
%\needspace{4\baselineskip}
\color{uniblue}{\subsection{Токовая защита ошиновки (ТЗО)}}
\color{black}

\begin{enumerate}[label=\arabic{section}.\arabic{subsection}.\arabic{enumi}, labelsep=4pt, leftmargin=0pt, itemindent=57pt, itemsep=0pt, parsep=0pt] % если что можно крутить itemsep 0 или 5

\item
Для ЛЭП, подключенной к РУ более чем через один выключатель, в устройстве предусмотрена токовая защита ошиновки (ТЗО), автоматически вводимая по факту отключенного положения линейного разъединителя (ЛР). Предусмотрена также возможность оперативного ввода (независимо от положения ЛР) от внешнего сигнала \textit{«Опер. ввод ТЗО»}.

\item
В составе ТЗО предусмотрены:
\begin{itemize}
    \item ступень, реагирующая на величину фазных токов (ТЗО-Ф);
    \item ступень, реагирующая на величину тока нулевой последовательности (ТЗО-НП);
\end{itemize}

\item
Ввод в работу ступеней осуществляется программными ключами \textbf{«ТЗО-Ф»} и \textbf{«ТЗО-НП»}.

\item
Токи срабатывания ТЗО-Ф и ТЗО-НП задаются с помощью уставок \textbf{«Iср ТЗО-Ф»} и \textbf{«3I0ср ТЗО-НП»} соответственно. Действие ступеней осуществляется с регулируемыми выдержками времени \textbf{«Tср ТЗО-Ф»} и \textbf{«Tср ТЗО-НП»}.

\item
ТЗО-Ф и ТЗО-НП могут контролироваться блокировкой при сквозных токах через ошиновку. Блокировка вводится отдельными программными ключами \textbf{«БСТО ТЗО-Ф»} и \textbf{«БСТО-НП»}.

\item
ТЗО может быть оперативно выведена из работы от внешнего сигнала \textit{«Вывод ТЗО»}, назначенного на соответствующий дискретный вход или функциональную клавишу.

\item
Функциональная схема логики ТЗО приведена на рисунке \ref{pict_tzo:TZO}.

\begin{figure}[H]
\centering
\includegraphics[width=1\textwidth,height=1\textheight,keepaspectratio]{\fbpath/561.0181 ТЗО/img/TZO.pdf}
\captionof{figure}{Функциональная схема ТЗО}\label{pict_tzo:TZO}
\end{figure}

\item
Уставки ТЗО представлены в таблице \ref{app_set:tzo}.

\end{enumerate}

%\newpage